\documentclass[11pt]{article}

% --- packages ---------------------------------------------------------------
\usepackage[utf8]{inputenc}
\usepackage[T1]{fontenc}
\usepackage[margin=1in]{geometry}
\usepackage{microtype}
\usepackage{amsmath}
\usepackage{booktabs}
\usepackage{graphicx}
\usepackage{xcolor}
\usepackage{listings}
\usepackage[numbers,sort&compress]{natbib}
\usepackage[colorlinks=true,allcolors=blue]{hyperref}

% --- listings (JSON / ascii diagrams) ---------------------------------------
\definecolor{codebg}{gray}{0.96}
\lstset{
  basicstyle=\ttfamily\footnotesize,
  backgroundcolor=\color{codebg},
  frame=single,
  framesep=4pt,
  rulecolor=\color{gray!40},
  breaklines=true,
  columns=fullflexible,
  keepspaces=true,
  showstringspaces=false,
  aboveskip=1em,
  belowskip=1em,
}

\newcommand{\code}[1]{\texttt{\small #1}}

% Include a screenshot if present, else a labeled placeholder box (keeps the
% build green before the PNGs are dropped into figures/).
\newcommand{\viewfig}[1]{%
  \IfFileExists{#1}{\includegraphics[width=\linewidth]{#1}}{%
    \fbox{\parbox[c][3cm][c]{0.92\linewidth}{\centering\footnotesize\itshape
      screenshot pending\\\texttt{\scriptsize #1}}}}}

% --- title ------------------------------------------------------------------
\title{\textbf{DSLs as Harnesses for Reliable LLM Agents}}

\author{
  Satoshi Nakajima\\
  \small The Singularity Society\\
  \small \texttt{satoshi.nakajima@gmail.com}
}

\date{\today}

\begin{document}
\maketitle

\begin{abstract}
A recurring lesson of applied agent engineering is that the \emph{harness}---the
designed environment an agent operates in---governs reliability more than the
underlying model. We argue that a \emph{domain-specific language} (DSL) is one of
the most effective harnesses available for autonomous, long-horizon agents,
precisely because it gives up the unbounded expressiveness of general-purpose
code in exchange for a small, legible, checkable surface. We make this concrete
by describing \textbf{Collections}, a declarative \emph{applications-as-data}
system implemented in MulmoClaude, an open-source LLM-agent application. In
Collections, a single JSON schema defines an application's data model,
relationships, computed values, recurrence, user interface, external-data
acquisition, and workflows; record files are the database; and the language model
is the runtime that operates inside the schema rather than writing imperative
code. We show how four well-known harness properties---schema-as-system-of-record,
a validator as a free feedback loop, a grammar as a forcing function, and a unit
of the language as a unit of context---fall out of this single design decision,
and how a deliberate \emph{reliability dial} divides what the host guarantees
deterministically (schema validity, derived-value computation, recurrence math,
sandboxed rendering, scheduled ingestion) from what the model supplies (judgment,
language, domain modeling). We position the approach against token-level
constrained decoding and schema-validated tool interfaces---which bound a single
generation or call---whereas a collection schema bounds a persistent, on-disk
application the agent operates \emph{inside}. We report a longitudinal single-user
deployment in which 24 schema-defined applications, 3 external feeds, and 9 bespoke
sandboxed UI views accreted over roughly two months without any host code being
written per application. We discuss the design's principal hazard---the escape hatch from the
declarative surface to arbitrary code---and its limitations: the open-ended
formula and custom-view sub-surfaces that lie outside the finite grammar, indirect
prompt injection through ingested feed content, non-replayable execution, and the
single-user, single-system scope of our evidence.
\end{abstract}

\noindent\textbf{Keywords:} LLM agents; agent harness; domain-specific languages;
end-user software engineering; low-code; applications as data; declarative systems.

% ===========================================================================
\section{Introduction}
\label{sec:intro}

The prevailing lesson of applied AI engineering in 2024--2026 is that the
\emph{harness}---the designed environment an agent operates inside: the tools it
can call, the format of the information it receives, the feedback that catches its
mistakes, and the scaffolding that lets it hand work to its future self---matters
more than the model alone~\citep{yang2024sweagent,yao2023react}. The same model,
given a better harness, produces dramatically better work. For a language model,
the interface is not a convenience layer wrapped around the intelligence; it is
the cognitive architecture.

This paper argues a narrower and, we believe, more actionable point: \textbf{a
domain-specific language (DSL) is one of the most powerful harnesses one can give
an agent, precisely because it gives up power.} A general-purpose programming
language is Turing-complete; that is its glory and, for an autonomous agent, its
hazard. There are infinitely many ways to express any behavior, no two codebases
agree on conventions, and nothing in the language tells the agent whether a given
program is \emph{correct for this domain}. Freedom under uncertainty is where
agents thrash: they explore broadly, drift toward locally plausible patterns, and
accumulate subtle errors that surface much later. A DSL inverts every one of those
properties---it expresses a bounded set of things, in a bounded number of ways,
and typically ships with a parser and validator that can decide well-formedness
before anything executes.

We ground the argument in a working system. \textbf{Collections} is a feature of
MulmoClaude, an open-source agent application built on the Claude Agent SDK whose
founding philosophy is \emph{the workspace is the database; files are the source
of truth; the model is the intelligent interface.} A collection is defined by a
single JSON \emph{schema} that simultaneously specifies a data model, inter-record
relationships, spreadsheet-like derived computations, completion and notification
rules, host-driven recurrence, a rendered user interface, optional bespoke HTML
views, and optional scheduled ingestion of external data. The records are plain
JSON files on disk. The host platform contains \emph{zero} domain knowledge; it
understands only generic concepts (fields, relationships, derived values, views,
actions, ingestion). Everything domain-specific lives in the schema, and the
language model operates the application by reading and writing records under the
schema's constraints.

The significance for reliability is that a self-built, model-operated application
need not be ``flaky chat with extra steps.'' The DSL is a contract that pins down
what must be deterministic while the model supplies flexibility and judgment, and
the two live on opposite sides of a deliberate line we call the \emph{reliability
dial} (\S\ref{sec:dial}). A schema \emph{validator} makes the boundary
self-enforcing: a structurally incoherent application is rejected at load, with an
actionable message, before it can run.

\paragraph{Contributions.}
\begin{enumerate}
  \item We articulate \textbf{DSL-as-harness}: a mapping from a single design
  decision (constrain the agent's authoring surface to a validated DSL) to four
  established harness properties, and an analysis of how that case \emph{shifts}
  rather than weakens as models improve (\S\ref{sec:harness}).
  \item We describe \textbf{Collections}, a concrete declarative
  applications-as-data system, including its schema DSL, relationships, derived
  formulas, host-driven recurrence, sandboxed custom views, and declarative feed
  ingestion (\S\ref{sec:collections}).
  \item We formalize the \textbf{host/model responsibility split} and the
  \emph{reliability dial}, and explain the validator and escape-hatch mechanisms
  that keep the surface both safe and expressive (\S\ref{sec:dial},
  \S\ref{sec:escape}).
  \item We report a \textbf{longitudinal single-user deployment}: 24 schema-defined
  applications, 3 external feeds, and 9 bespoke sandboxed views accreted over
  $\sim$2 months with no per-application host code (\S\ref{sec:casestudy}).
  \item We give an \textbf{honest account of the limitations}: the escape-hatch
  hazard, non-replayable execution, and the scope of our evidence
  (\S\ref{sec:limitations}).
\end{enumerate}

% ===========================================================================
\section{Background and Related Work}
\label{sec:related}

\paragraph{Agent harnesses and agent--computer interfaces.}
\citet{yang2024sweagent} showed that a purpose-built agent--computer interface
(ACI)---structured tools, capped and guided search, and a linter that rejects
syntactically broken edits at the moment of introduction---improves software-engineering
agent performance substantially with a fixed model. \citet{yao2023react}
interleave reasoning and acting; \citet{schick2023toolformer} teach models to call
external tools. These works establish the empirical primacy of the surrounding
environment. Our contribution is to identify the \emph{DSL} as a particularly
sharp instance of such an environment and to instantiate it as an
applications-as-data substrate.

\paragraph{Self-improving and memory-augmented agents.}
Voyager~\citep{wang2023voyager} accretes a reusable skill library, and Generative
Agents~\citep{park2023generative} maintain a memory stream they reflect over. Both
let an agent extend itself over time. Collections differ in substrate: the agent's
accreted capabilities are \emph{declarative DSL artifacts on disk}, not learned
weights or imperative code, which makes them inspectable, diffable, and validated.
A complementary line of work makes the event log the primary
artifact---\citet{ynakajima2026log} prioritizes an append-only log to obtain
deterministic replay and cheap forking. We deliberately forgo byte-for-byte replay
(the model step is nondeterministic) in exchange for an agent whose configuration
is plain, readable, validated files; \S\ref{sec:limitations} returns to this
trade.

\paragraph{Domain-specific languages.}
The software-engineering literature has long studied when and how a deliberately
limited language pays off~\citep{mernik2005dsl,fowler2010dsl}: a DSL trades general
expressiveness for analyzability, a smaller error surface, and tooling (parsers,
validators) that come ``for free.'' We apply this classical argument to a new
consumer---an autonomous language-model agent---and observe that the properties
DSLs were always praised for map directly onto modern harness-engineering
guidance.

\paragraph{Constrained decoding and structured outputs.}
A parallel line of work enforces structural validity at \emph{inference} time:
grammar- and schema-constrained decoding masks the token distribution so that only
well-formed continuations are sampled (Outlines~\citep{willard2023outlines}, LMQL~\citep{beurerkellner2023lmql},
Guidance), and OpenAI's function calling and Structured
Outputs~\citep{openai2024structured} make a JSON schema govern a model's output for
a single call. These mechanisms realize our property
(3)---grammar as a forcing function---at the \emph{token/model boundary}, and they
do so more tightly than we do: Collections uses \emph{post-hoc} JSON-schema
validation plus retry, enforcing well-formedness at the \emph{system boundary}
rather than during generation. The distinction we draw is one of \emph{scope and
persistence}, not of mechanism: constrained decoding shapes one generation, and a
function/Structured-Outputs schema shapes one call's arguments, whereas a
collection schema is a durable, on-disk definition of an entire application's data
model, lifecycle, and UI that the agent operates \emph{inside} over many turns. We
adopt the inference-time techniques as complementary---they could tighten how the
agent emits a schema---rather than as alternatives to the application-as-data
substrate.

\paragraph{Schema-validated agent frameworks.}
Agent frameworks already place machine-checkable schemas between an agent and its
environment: LangChain's typed (Pydantic) tool definitions, AutoGen's agent
configuration~\citep{wu2023autogen}, and Semantic Kernel's plugin manifests all
validate the structure of \emph{individual tool calls}. This is properties (1) and
(2) at the granularity of a callable interface. Collections differs in where the
schema sits: rather than defining the boundary of a call the agent makes
\emph{across}, the schema defines the interior structure of an application the
agent inhabits---a persistent, evolvable configuration the host renders and
reconciles, not a per-invocation contract. The novelty we claim is therefore not
``schema-validated agent interface'' (which is established prior art) but
``schema-\emph{as-the-application}, with the agent as its runtime.''

\paragraph{End-user software engineering and low-code.}
End-user programming has sought for decades to let non-programmers build software,
from the spreadsheet onward~\citep{nardi1993small,ko2011enduser}; commercial
low-code/no-code platforms (Airtable, Notion, Retool, PowerApps) continue this
lineage, and by 2024--2026 several added AI-assisted, intent-based authoring
(Airtable's field generation, Notion AI's schema construction from natural
language, Power Automate's Copilot)---so ``intent-driven authoring'' is no longer
unique to our setting. The persistent constraint in that tradition is nonetheless
that the schema remains a \emph{UI/data-layer supplement} a human owns and the AI
\emph{assists}, with validation optional. Collections differs on three axes: the
schema is the \emph{complete} application definition (data model, computation,
recurrence, UI, ingestion) rather than a layer of one; validation is
\emph{architectural} (a structurally incoherent schema cannot load) rather than
advisory; and the agent \emph{authors the whole DSL program} from intent rather
than assisting a human composer. We frame this as the shift from \emph{no-code}
(a human composes blocks) to \emph{intent-driven} (an agent composes, and
operates, the DSL program).

% ===========================================================================
\section{A DSL Is a Harness}
\label{sec:harness}

When you hand an agent a DSL you are not merely giving it a tool; you are reshaping
the space it must search. That reshaping makes four classic harness patterns fall
out of a single design decision.

\paragraph{(1) The schema is the system of record.}
Harness engineering insists that anything the agent cannot read in context
effectively does not exist; intent must live in a machine-readable artifact. A DSL
satisfies this by construction: its schema \emph{is} the specification. When the
agent writes a DSL document, the user's intent is captured as a structured,
inspectable object rather than dissolved into the incidental details of imperative
code. It can be diffed, versioned, validated, and reasoned about without running
it.

\paragraph{(2) The validator is a free feedback loop.}
The highest-leverage component in the SWE-agent study was a linter that rejected
broken edits at the moment of introduction, before the error could propagate
through later steps~\citep{yang2024sweagent}. Every DSL ships with this loop:
the instant the agent emits a malformed structure, a parser or schema checker
rejects it with a localized, actionable message. The feedback loop came bundled
with the language.

\paragraph{(3) The grammar is a forcing function.}
Because the vocabulary is finite, the room to hallucinate is structurally smaller.
The agent cannot invent an option absent from the grammar; it can only compose the
primitives the language provides. Constraint, here, is a feature. Two
qualifications matter. First, Collections enforces this property
\emph{post hoc}---a parser/validator rejects a malformed schema and the agent
retries---rather than at the token level as constrained decoding does
(\S\ref{sec:related}); the forcing function therefore acts at the system boundary,
not inside the generation step. Second, and more importantly, the finite-grammar
guarantee covers the schema's \emph{structural} vocabulary---field types,
relationship kinds, layout and recurrence declarations---but \emph{not} the two
open-ended sub-surfaces the schema may embed: derived-field formula strings and
custom-view HTML. We bound the formula sub-language deliberately (it is a small,
non-Turing-complete expression grammar; \S\ref{sec:collections}), but a
\emph{well-formed} formula or view can still be semantically wrong, which we treat
as a residual reliability surface in \S\ref{sec:limitations}.

\paragraph{(4) The unit of the language is the unit of context.}
A well-designed DSL has a natural decomposition---a field, a record, a rule---which
becomes the agent's editing unit. It can work on one piece without holding the
whole document in mind, and the language's structure, not an ad-hoc heuristic,
defines the seams (progressive disclosure).

\paragraph{How the value shifts as models improve.}
A tempting intuition holds that DSLs are a crutch for weak models, to be discarded
as models get stronger. We argue the value does not vanish but \emph{shifts among
the four properties}, and we are careful to separate two claims the objection
conflates. The strongest form of the objection is about \emph{capability scaling}:
a frontier model emits structurally valid schemas at near-unit frequency, so the
marginal benefit of properties (2) the validator and (3) the grammar---both of
which catch \emph{structural} errors---does indeed tend toward zero as structural
generation improves. We concede this. But properties (1) and (4) are not error-correction:
making intent a diffable, inspectable artifact (1) and decomposing it into editing
units (4) are valuable independent of how good the model is at avoiding mistakes;
they govern \emph{auditability and progressive disclosure}, which a more capable
model exploits more, not less. Separately, there is a \emph{deployment-context}
claim: when work is delegated autonomously, over long horizons, with no human
reviewing each step and a requirement that execution be verifiable and reversible,
a persistent validated contract remains valuable regardless of model strength---not
to catch the model's structural mistakes but to make the system's behavior
checkable by something other than trust. A stronger model inside a DSL harness is
thus not wasted capability; it is capability aimed at judgment---what to model,
what to express---while the language supplies durability and inspectability that
scaling does not provide on its own.

% ===========================================================================
\section{The Collections System}
\label{sec:collections}

Collections applies the DSL-as-harness idea to software itself. A collection is:
\begin{center}
\code{schema.json} \;+\; \code{records/*.json} \;+\; model \;=\; application.
\end{center}
There is no database server, no migration framework, and no application-specific
backend. The host understands only generic capabilities; everything
domain-specific lives in the collection.

\subsection{The schema DSL}
A schema declares typed \emph{fields}. The current implementation supports 17
field types, including \code{string}, \code{text}, \code{number}, \code{date},
\code{datetime}, \code{boolean}, \code{markdown}, \code{enum}, \code{money},
\code{image}, \code{file}, a nested \code{table} of sub-records, two relationship
types (\code{ref}, \code{embed}), a computed \code{derived} field, and a
\code{toggle} that projects an enum. Relationships are first-class: a \code{ref}
stores a foreign key and the host renders pickers, links, and navigation
generically; an \code{embed} pulls a fixed record from another collection for
display.

\subsection{Computation without code}
A \code{derived} field behaves like a spreadsheet formula, but formulas may follow
references into other collections:
\begin{lstlisting}
"total": {
  "type": "derived", "display": "money", "currencyField": "currency",
  "formula": "sum(lineItems[].quantity * lineItems[].rate) * (1 + taxRate)"
}
\end{lstlisting}
The host evaluates derived fields, and crucially returns them \emph{to the model}
as well as to the screen, so the agent reasons over the same numbers the user
sees. A portfolio can revalue itself when a quote changes elsewhere, with no
synchronization code.

The formula is itself a small DSL, and a deliberately bounded one. Its grammar
admits only the four arithmetic operators, parentheses, a single aggregation
function (\code{sum} over a table column or product of table columns), and
single-level reference dereferencing (e.g.\ \code{ticker.price}); it has no
conditionals, loops, string operations, comparisons, user-defined functions, or
nested calls. It is therefore \emph{not} Turing-complete---closer to a spreadsheet
cell expression than a scripting language---and it is parsed before evaluation, so
a malformed formula, an unbound identifier, or a non-finite result yields a null
that the host renders as an em-dash rather than a wrong number. This bound limits
the \emph{structural} surface of a formula; it does not make a well-formed formula
\emph{semantically} correct, a gap we return to in \S\ref{sec:limitations}.

\subsection{Completion, notification, and host-driven recurrence}
A schema can mark which field and values count as ``done''
(\code{completionField} / \code{completionDoneValues}), gate notifications on a
date (\code{triggerField} / \code{triggerLeadDays}), and---via a \code{spawn}
block---provision the next instance of a recurring obligation on a civil-date
cadence when the current one closes:
\begin{lstlisting}
"triggerField": "dueOn", "triggerLeadDays": 3,
"spawn": {
  "every": { "unit": "month", "interval": 1, "dayOfMonth": 10 },
  "carry": ["payee", "amount"], "set": { "status": "pending" }
}
\end{lstlisting}
The author declares \emph{what recurs and when to nudge}; the host's reconciler
owns the deterministic half---raising and clearing reminders and advancing each
cycle with create-if-absent (idempotent) semantics. No timer, cron expression, or
notification code is written.

\subsection{Custom views: bespoke UI as data}
Host-rendered field types give every collection a list, a detail form, and---when
declared---calendar and Kanban layouts for free. When field types cannot present
the data well (a portfolio as an allocation chart, a set of places as a map, a
word list as flashcards), a collection may carry a \emph{custom view}: an HTML
file the agent authors and the host serves.
\begin{lstlisting}
"views": [
  { "id": "allocation", "label": "Allocation",
    "file": "views/allocation.html", "capabilities": ["read", "write"] }
]
\end{lstlisting}
A view is an ordinary web page (it may load a charting library from a curated CDN
allow-list) but it is rendered inside a sandboxed \code{iframe} under a strict
Content-Security-Policy: an opaque origin, \code{default-src 'none'}, and network
access only back to the host. It receives its records through a capability-scoped,
time-limited token rather than ambient access; the declared \code{capabilities}
(\code{read}, \code{write}) bound what that token permits. A custom view is thus
not a hole in the harness but the same bargain as the rest of the system: the
model supplies the bespoke surface, the host supplies the deterministic, sandboxed
boundary, and the presentation layer too becomes inspectable, versioned data.

\subsection{Feeds: acquisition as data}
Some applications track an external source---a podcast's episodes, a newspaper's
headlines, a JSON API. Such a collection declares an \code{ingest} block and
becomes a \emph{feed}:
\begin{lstlisting}
"ingest": {
  "kind": "rss", "url": "https://lexfridman.com/feed/podcast/",
  "schedule": "daily", "idFrom": "guid",
  "map": { "headline": "title", "url": "link", "published": "pubDate" },
  "maxItems": 100
}
\end{lstlisting}
\code{kind} selects a retriever (\code{rss}, \code{atom}, \code{http-json});
\code{map} projects source fields onto collection fields; \code{schedule}
(\code{hourly}/\code{daily}/\code{weekly}/\code{on-demand}) sets cadence;
\code{maxItems} caps retention. The host owns fetching, parsing, field mapping,
de-duplication by id, pruning, and cursor/failure tracking. Because a feed is
still a collection, the user can layer their own fields onto the ingested
ones---a rating, a note, a resume position---turning a passive mirror into a
read--write application over an external source.

\subsection{Schema evolution}
\label{sec:evolution}
Because an agent edits schemas in place as needs change, records written under an
earlier schema coexist with a later one, and the host must define what happens.
Its rule is asymmetric: schemas are validated at \emph{load} (the path-independent
gate of \S\ref{sec:dial}) and records are validated on \emph{write}, but records
are \emph{not} re-validated on \emph{read}. The consequence is that
\emph{additive} changes are non-breaking by construction: a record lacking a
newly-added field is read without error, the missing value renders as empty, and
any derived field that referenced it evaluates to the em-dash null rather than
failing. A newly-\emph{required} field or a field-type change is breaking only
\emph{prospectively}---existing records still load and display, but the next write
to such a record (or a new record) must satisfy the new constraint. There is no
migration framework, and this is deliberate: re-validating on read would make a
single stale record look like data loss, so the host favors graceful degradation
over hard failure, and leaves true migrations (backfilling a required field across
old records) to an ordinary agent task. This behavior is exercised in our
deployment, not hypothetical: the podcast feed of \S\ref{sec:casestudy} gained
user rating and notes fields after ingestion had populated it, and the pre-existing
records absorbed the additive change without migration.

% ===========================================================================
\section{The Reliability Dial}
\label{sec:dial}

The reason a self-built, model-operated application is reliable rather than flaky
is a single principle: \textbf{the DSL is a contract that pins down what must be
deterministic, while the model supplies flexibility, judgment, and language.}
Table~\ref{tab:split} summarizes the split.

\begin{table}[t]
\centering
\small
\begin{tabular}{@{}ll@{}}
\toprule
\textbf{Host guarantees (deterministic)} & \textbf{Model owns (judgment)} \\
\midrule
Schema validity (load-time)        & What to model \\
Path safety, record storage        & What a record means \\
Derived-value computation          & How a workflow should proceed \\
Recurrence / civil-date math       & When to hand off to a human \\
Sandboxed view rendering (CSP)     & How to present data (authoring the view) \\
Scheduled ingestion + de-dup       & Which external source, mapped how \\
Record-write validation (every path) & Conducting a seeded-chat workflow \\
Load-time schema re-validation     & Which fields/views to add as needs evolve \\
\bottomrule
\end{tabular}
\caption{The host/model responsibility split. Everything on the left runs
identically every time; everything on the right is where the model's flexibility
is wanted.}
\label{tab:split}
\end{table}

A schema validator makes the boundary self-enforcing. A structurally incoherent
application is rejected at load with an actionable reason---examples include a
money field with neither a literal currency nor a per-record currency field, a
\code{ref} pointing at a non-existent collection, a \code{triggerField} that does
not name a real date field, and a \code{spawn} whose successor would be born
already matching its own predicate (an unbounded respawn).

The craft is \emph{where the dial sits}. Push everything into the DSL and you
rebuild the brittle ``can't-express-this'' cage that made no-code platforms
frustrating; push everything to the model and you lose the durable contract that
made it an application at all. The guiding rule is to \emph{extend the declarative
layer only when it outperforms the agent}.

% ===========================================================================
\section{The Escape Hatch}
\label{sec:escape}

The same rigidity that makes a DSL a good harness is its failure mode: when the
language cannot express what the user needs, an agent will either give up or
contort the DSL into something it was never meant to hold. Every well-designed DSL
answers this with a deliberate exit to a more expressive layer.

Collections provides two. First, a per-record or collection-level \code{action}
hands off to a seeded chat---a full-tool agent---exactly when a task needs human
judgment the form cannot capture (``generate the invoice PDF,'' ``what was the
amount, and was it actually paid?''). This is business logic as prose, behind a
declarative door. Second, the \emph{custom view} (\S\ref{sec:collections}) is the
escape for presentation: when host-rendered field types cannot express a layout,
the schema drops to arbitrary HTML in a sandbox rather than distorting the
field-type vocabulary. The art is not to maximize constraint but to choose
\emph{where} the constraint binds and to provide a clean, legible path out where it
does not. A harness with no escape hatch eventually forces the agent to fail or to
lie, and both are worse than a slightly leakier abstraction.

It is important to be precise about how bounded the seeded-chat escape actually is,
because ``full-tool agent'' overstates it. The chat runs under a named \emph{role},
and the role's allow-list determines the tool surface; the agent in the seeded chat
can do only what that role grants. Record reads and writes go through the same
\emph{record-write validation} that governs every other path, so a seeded chat
cannot write a record the schema would reject. Crucially, there is no tool that
edits a schema directly, and even if a schema file were modified by any path, it is
re-validated by the same load-time validator before the collection can be used
again (\S\ref{sec:dial}): the validator is a \emph{path-independent gate}, not a
write-time check that an escape could sidestep. What the escape genuinely relaxes is
the \emph{conduct of the workflow}---the conversation itself is model-driven and
unconstrained---not the structural guarantees over what gets persisted. The
honest residual, then, is semantic rather than structural, which the next section
takes up.

% ===========================================================================
\section{Case Study: A Longitudinal Single-User Deployment}
\label{sec:casestudy}

We report on one author's production workspace, which has accreted applications
through ordinary use over roughly two months. No host code was written for any of
these applications; each is a schema plus records, authored by the agent in
response to natural-language requests.

\paragraph{Scale.}
The workspace contains 24 schema-defined collections holding $\sim$700 records,
plus 3 feeds (\S\ref{sec:collections}) holding a further $\sim$260 ingested
records. Collection sizes vary by nearly three orders of magnitude
(Table~\ref{tab:collections}), consistent with the claim that the marginal cost of
a new application is low enough to justify even a single-record one. Each schema is
small---32 lines of JSON on average---and there is, by construction, zero
per-application host code behind any of them.

\begin{table}[t]
\centering
\small
\begin{tabular}{@{}lr lr lr@{}}
\toprule
Collection & \#rec & Collection & \#rec & Collection & \#rec \\
\midrule
mag2-subscribers & 466 & engagements & 22 & recurring-bills & 6 \\
vocabulary       & 50  & restaurants & 20 & videos          & 6 \\
todos            & 23  & lessons-biology & 12 & worklog      & 6 \\
reading-list     & 10  & baseball-scout  & 11 & portfolio    & 4 \\
stock-quotes     & 10  & watchlist       & 10 & html-apps    & 4 \\
tax-payments     & 7   & clients         & 3  & \emph{(7 more)} & $\le 2$ \\
\bottomrule
\end{tabular}
\caption{Representative collections by record count in the studied workspace
(24 schema-defined collections total). Sizes span from a single-record profile to
466 newsletter subscribers.}
\label{tab:collections}
\end{table}

\paragraph{Mechanisms exercised in practice.}
The deployment uses every major mechanism described in \S\ref{sec:collections}:
relationships (an \code{invoice} references a \code{client} and embeds the issuer
profile), derived computation (invoice totals from line items and tax),
host-driven recurrence (monthly and annual obligations via \code{spawn}), and
completion/notification gating (todos, tax payments).

\paragraph{Bespoke views.}
Nine custom HTML views, authored by the agent on request, are deployed---seven over
collections and two over feeds. They include a restaurant guide rendered as a
transit-style map, a movie watchlist as a poster gallery, a vocabulary list as a
flashcard drill, a portfolio as an allocation chart, a stock list as a valuation
scatter plot, a scouting collection as per-entity radar charts, and a subscriber
collection as a trend line (Figure~\ref{fig:views}). Each is software with an
addressable audience of one---software no vendor would build---produced for the cost
of a request. The owner never wrote or read the HTML.

\begin{figure}[t]
\centering
\begin{minipage}[t]{0.31\textwidth}\centering
  \viewfig{figures/restaurants-map.png}\\[2pt]\footnotesize (a) a restaurant guide as a transit-style map
\end{minipage}\hfill
\begin{minipage}[t]{0.31\textwidth}\centering
  \viewfig{figures/portfolio-allocation.png}\\[2pt]\footnotesize (b) a portfolio as an allocation chart
\end{minipage}\hfill
\begin{minipage}[t]{0.31\textwidth}\centering
  \viewfig{figures/feed-gallery.png}\\[2pt]\footnotesize (c) a podcast feed as a read--write gallery
\end{minipage}
\caption{Three of the nine agent-authored custom views in the studied deployment
(two collection views and one feed view): a restaurant collection rendered as a
transit-style map (a); a portfolio whose host-computed valuations (\S\ref{sec:collections})
drive an allocation chart (b); and an ingested podcast feed turned into a
read--write gallery (c), where the owner's rating, notes, and resume position are
layered on top of episode metadata. Each is a self-contained HTML file the host
serves in a sandboxed, capability-scoped iframe; the owner described each in
natural language and never saw the generated markup.}
\label{fig:views}
\end{figure}

\paragraph{Feeds.}
Three collections are feeds: an hourly newspaper-technology feed (64 records at
time of writing), and two daily podcast feeds (100 records each). One podcast feed
is read--write: the schema layers the owner's own rating, notes, and
resume-position fields on top of the ingested episode metadata, and the host
serves a custom gallery view over the combined record.

\paragraph{Observations.}
Two qualitative observations stand out. (i) \emph{No per-application host code} was
required for any of the 26 applications; the host's generic capabilities sufficed,
which is the load-bearing claim of applications-as-data. This includes the one
instance of schema evolution we observed---the podcast feed gaining user rating and
notes fields after ingestion---which the host absorbed as an additive,
non-breaking change with no migration code (\S\ref{sec:evolution}); we did not
stress breaking changes (a newly-required field or a type change across a large
record set), so the no-migration claim is bounded to additive evolution. (ii) The applications
\emph{diverge sharply from any mass-market product}: a transit-map restaurant
guide and a read--write podcast tracker with personal resume positions reflect one
user's specific needs, supporting the view that a near-zero marginal build cost
shifts the natural unit of software from the mass-market product toward the
individual.

% ===========================================================================
\section{Discussion and Limitations}
\label{sec:limitations}

\paragraph{The escape-hatch hazard.}
The expressiveness that escape hatches restore is also unconstrained: a custom
view or a seeded-chat action runs outside the declarative guarantees. We mitigate
this for views with sandboxing and capability-scoped tokens (\S\ref{sec:collections}),
but the general tension---more expressiveness, fewer guarantees---is fundamental
and must be managed per escape, not eliminated.

\paragraph{Non-replayable execution.}
Because the operating step is a nondeterministic model call, a run cannot be
replayed byte-for-byte, and a self-modification's effect cannot be A/B-replayed
against the prior behavior on identical inputs. This is the deliberate trade
against log-primary designs~\citep{ynakajima2026log}. The mitigation is
configuration-level rather than run-level: a git-backed workspace makes every
schema and view edit auditable and revertable, which is the granularity that
matters for an agent editing its own applications.

\paragraph{Unconstrained sub-surfaces within the schema.}
As noted in \S\ref{sec:harness} and \S\ref{sec:collections}, the finite-grammar
guarantee covers the schema's structural vocabulary but not the two expressive
sub-surfaces it can embed: derived-field formula strings and custom-view HTML.
These are where the harness's structural guarantee ends and ordinary
model-reliability concerns resume. We bound the formula grammar (non-Turing-complete,
parsed before evaluation) and sandbox the views (opaque-origin iframe, strict CSP,
capability-scoped tokens), so the \emph{blast radius} of an error in either is
contained; what is \emph{not} contained is whether a well-formed formula or a
well-rendered view is \emph{correct}. Structural validity and semantic correctness
are different guarantees, and the DSL provides only the first.

\paragraph{Correctness has no second audience.}
Mass software accumulates correctness through many users and years; an application
with an audience of one does not. A wrong poster gallery is harmless, but a wrong
financial or tax view is not, and only the model and the single user ever inspect
it. The reliability dial is the relevant defense---push consequential computation
into host-guaranteed derived fields rather than leaving it to a generated
view---but it \emph{relocates} the valid-but-wrong failure mode rather than
removing it. A derived formula is itself a model-authored string: the host
evaluates it deterministically but cannot judge whether it is the right formula for
the domain. A plausible-but-erroneous tax formula---say, applying the rate before a
discount rather than after---passes schema validation and then executes reliably
and wrongly on every record, with no second reviewer to catch it. The mitigation
that the substrate does afford is that formulas are short, inspectable, plain-text
expressions in the schema file, so \emph{diffing a formula against the user's
stated intent} is a tractable review step in a way that auditing generated
imperative code is not; but the residual risk is real and is part of the price of
building for one.

\paragraph{Indirect prompt injection via ingested content.}
Feeds (\S\ref{sec:collections}) introduce a threat the sandboxing does not address.
Ingested fields---headlines, descriptions, show notes---are untrusted natural
language that then enters records the model reads and reasons over, and indirect
prompt injection (adversarial instructions embedded in retrieved content that the
model subsequently follows) is a documented attack on content-consuming
agents~\citep{greshake2023injection}. The custom-view CSP sandbox governs UI
rendering, not the model's context, so it offers no protection here; and because
the operating agent may hold write-capable tokens, a compromised feed could in
principle steer record writes or cross-collection reads within whatever permissions
the active role grants. Capability scoping bounds \emph{what} an injected
instruction could reach but does not stop the model from acting on embedded text in
the first place. We did not observe such an attack in our benign deployment, but we
flag it as the feed architecture's primary unmitigated threat; plausible defenses
include processing ingested content in a system-prompt-isolated context and
flagging model operations initiated over recently-ingested records for confirmation.

\paragraph{Threats to validity.}
Our empirical evidence is a single user, a single system, and an author-operated
deployment; it demonstrates feasibility and the no-host-code property but does not
establish generality, and it reports usage rather than a controlled comparison.
Quantifying authoring effort, error rates, and how the reliability dial should be
set across users and domains is future work.

% ===========================================================================
\section{Conclusion}
\label{sec:conclusion}

A domain-specific language is a powerful harness for an autonomous agent precisely
because it relinquishes expressiveness for a small, validated, inspectable
surface. Collections applies that principle to applications themselves: a schema
defines the application, record files are the database, the host runs the
deterministic half, and the model operates inside the schema. The result is
software an agent can author and refine on request, with the reliability of a
declarative contract and the flexibility of a language model held apart by a
deliberate dial. Building applications shifts from writing imperative code to
declaring---and validating---harnesses; the records are data, the schema is the
program, and the model is the runtime.

% ===========================================================================
\section*{Availability}
MulmoClaude and the Collections implementation are open source. The schema
validator, derived-formula evaluator, recurrence reconciler, custom-view sandbox,
and feed engine referenced above are part of that codebase, available at
\url{https://github.com/receptron/mulmoclaude}.

\bibliographystyle{plainnat}
\bibliography{refs}

\end{document}
